\documentclass[11pt]{article}
\usepackage[utf8]{inputenc}
\usepackage[T1]{fontenc}
\usepackage[spanish]{babel}
\usepackage{amsmath, amssymb, bm}
\usepackage{geometry}
\geometry{margin=2.5cm}

\title{Ridge/LS 04: eigenvalues, SVD y estabilidad numerica}
\author{}
\date{}

\begin{document}
\maketitle

\section{Objetivo}

Este archivo explica que ocurre con:
\[
    (A^H A + \lambda I)^{-1}A^H
\]
desde la perspectiva espectral. Es decir, usando eigenvalues, singular values y
direcciones principales.

\section{Descomposicion en valores singulares}

Toda matriz compleja \(A\in\mathbb{C}^{M\times L}\) puede escribirse como:
\[
    A = U\Sigma V^H.
\]

Aqui:

\begin{itemize}
    \item \(U\) contiene direcciones ortonormales en el espacio de mediciones.
    \item \(V\) contiene direcciones ortonormales en el espacio de parametros.
    \item \(\Sigma\) contiene los valores singulares \(\sigma_i\).
\end{itemize}

Los valores singulares cumplen:
\[
    \sigma_i \ge 0.
\]

\section{Relacion con \(A^H A\)}

Si:
\[
    A = U\Sigma V^H,
\]
entonces:
\[
\begin{aligned}
    A^H A
    &=
    (U\Sigma V^H)^H(U\Sigma V^H)\\
    &=
    V\Sigma^H U^H U\Sigma V^H\\
    &=
    V\Sigma^H\Sigma V^H.
\end{aligned}
\]

Como \(U^H U = I\), queda:
\[
    A^H A = V\Sigma^2 V^H.
\]

Por tanto, los eigenvalues de \(A^H A\) son:
\[
    \mu_i = \sigma_i^2.
\]

\section{LS puro visto con SVD}

El estimador LS puro es:
\[
    \hat{\bm h}_{LS} = (A^H A)^{-1}A^H\bm y.
\]

Usando SVD:
\[
    A^\dagger = V\Sigma^\dagger U^H.
\]

Esto significa:
\[
    \hat{\bm h}_{LS}
    =
    \sum_i
    \frac{1}{\sigma_i}
    \bm v_i
    \bm u_i^H
    \bm y.
\]

Cada componente de \(\bm y\) en la direccion \(\bm u_i\) se amplifica por:
\[
    \frac{1}{\sigma_i}.
\]

Si \(\sigma_i\) es pequeno, esa direccion se amplifica muchisimo.

\section{Riesgo de valores singulares pequenos}

Supongamos:
\[
    \bm y = A\bm h + \bm n.
\]

Entonces:
\[
    \hat{\bm h}_{LS}
    =
    A^\dagger A\bm h
    +
    A^\dagger \bm n.
\]

El segundo termino es ruido filtrado:
\[
    A^\dagger \bm n.
\]

Si algun \(\sigma_i\) es pequeno, el ruido en esa direccion se multiplica por:
\[
    \frac{1}{\sigma_i}.
\]

Este es uno de los motivos principales para usar regularizacion.

\section{Ridge visto con SVD}

El estimador ridge es:
\[
    \hat{\bm h}_{ridge}
    =
    (A^H A+\lambda I)^{-1}A^H\bm y.
\]

Usando:
\[
    A^H A = V\Sigma^2 V^H,
\]
tenemos:
\[
    A^H A+\lambda I
    =
    V(\Sigma^2+\lambda I)V^H.
\]

Su inversa es:
\[
    (A^H A+\lambda I)^{-1}
    =
    V(\Sigma^2+\lambda I)^{-1}V^H.
\]

Ademas:
\[
    A^H = V\Sigma^H U^H.
\]

Entonces:
\[
    (A^H A+\lambda I)^{-1}A^H
    =
    V(\Sigma^2+\lambda I)^{-1}\Sigma^H U^H.
\]

\section{Factores de filtrado}

El estimador ridge se puede escribir como:
\[
    \hat{\bm h}_{ridge}
    =
    \sum_i
    \frac{\sigma_i}{\sigma_i^2+\lambda}
    \bm v_i
    \bm u_i^H
    \bm y.
\]

El factor:
\[
    f_i(\lambda)
    =
    \frac{\sigma_i}{\sigma_i^2+\lambda}
\]
reemplaza al factor LS puro:
\[
    f_i(0) = \frac{1}{\sigma_i}.
\]

\section{Que pasa si \(\sigma_i\) es pequeno}

En LS puro:
\[
    f_i(0)=\frac{1}{\sigma_i}.
\]

Si \(\sigma_i\to 0\), entonces:
\[
    f_i(0)\to\infty.
\]

En ridge:
\[
    f_i(\lambda)
    =
    \frac{\sigma_i}{\sigma_i^2+\lambda}.
\]

Si \(\sigma_i\to 0\), entonces:
\[
    f_i(\lambda)\to 0.
\]

Esto significa que ridge no amplifica direcciones poco confiables; las atenua.

\section{Interpretacion como filtro}

Ridge actua como un filtro espectral:

\begin{itemize}
    \item Direcciones con \(\sigma_i\) grande: se conservan.
    \item Direcciones con \(\sigma_i\) pequeno: se reducen.
\end{itemize}

Por eso regularizar mejora estabilidad cuando el problema esta mal condicionado.

\section{Eigenvalues desplazados}

Como:
\[
    \mu_i = \sigma_i^2,
\]
la matriz regularizada tiene eigenvalues:
\[
    \mu_i+\lambda
    =
    \sigma_i^2+\lambda.
\]

Por tanto:
\[
    \lambda
\]
pone un piso minimo a los eigenvalues.

\section{Numero de condicion}

Para \(A^H A\):
\[
    \kappa(A^H A)
    =
    \frac{\sigma_{\max}^2}{\sigma_{\min}^2}.
\]

Para la matriz regularizada:
\[
    \kappa(A^H A+\lambda I)
    =
    \frac{\sigma_{\max}^2+\lambda}{\sigma_{\min}^2+\lambda}.
\]

Si \(\sigma_{\min}^2\) es muy pequeno, la regularizacion reduce mucho el numero
de condicion.

\section{Bias y varianza}

Ridge introduce un compromiso:

\begin{itemize}
    \item Reduce varianza: amplifica menos el ruido.
    \item Introduce bias: ya no ajusta exactamente como LS puro.
\end{itemize}

Ese compromiso suele ser deseable cuando hay ruido.

En estimacion de canal, esto significa:

\begin{itemize}
    \item LS puro puede seguir demasiado el ruido de los pilotos.
    \item Ridge produce un canal estimado mas suave y estable.
\end{itemize}

\section{Conexion con el simulador}

En el programa, \(A\) es una matriz parcial de Fourier. Sus valores singulares
dependen de:

\begin{itemize}
    \item posiciones de pilotos;
    \item cantidad de pilotos;
    \item cantidad de taps estimados;
    \item patron escalonado entre bloques.
\end{itemize}

La regularizacion:
\[
    \lambda = 10^{-2}
\]
evita que direcciones mal observadas del canal generen taps demasiado grandes.

\section{Conclusion}

La forma:
\[
    (A^H A+\lambda I)^{-1}A^H
\]
no solo es una formula algebraica. Es un filtro espectral:
\[
    \frac{\sigma_i}{\sigma_i^2+\lambda}.
\]

Esa es la clave para entender por que es mas estable que LS puro.

\end{document}
